\setcounter{table}{0}\setcounter{figure}{0}\setcounter{equation}{0}
\renewcommand{\thetable}{A.\arabic{table}}\renewcommand{\thefigure}{A.\arabic{figure}}\renewcommand{\theequation}{A.\arabic{equation}}

\section{Implementation checks}
\label{app:impl}

\paragraph{Augmentation on the GPU}
Every transform draws its random parameters per image, as \texttt{torchvision} does, but processes a whole
batch at once on the GPU. With fixed parameters we compared each operation with
\texttt{torchvision.transforms.functional} on 64 random images: brightness, saturation, hue and grayscale
match exactly; contrast to $6\times10^{-8}$; resized crop (bilinear) and crop with flip to
$1.2\times10^{-7}$. The test is part of the repository.

\paragraph{Loss}
\Cref{eq:loss} was compared with the official \supcon implementation on random embeddings. The maximum
absolute difference is $0$ when all images are labeled (\supcon), $4.8\times10^{-7}$ when none is (NT-Xent),
and $9.5\times10^{-7}$ when half of the batch is labeled (\union, compared with the official code and an
explicit mask).

\needspace{0.45\textheight}
\section{Hyper-parameters}
\label{app:hparams}
\Cref{tab:hparams} lists the hyper-parameters of every method and evaluation protocol.

\begin{table}[!h]
\centering
\small
\caption{Hyper-parameters of all methods (ResNet-18 on CIFAR unless stated). One contrastive step processes
1,024 augmented views.}
\label{tab:hparams}
\begin{adjustbox}{max width=\linewidth}
\begin{tabular}{l l}
\toprule
\simclr, \supcon, \union & SGD, momentum 0.9, lr 0.25 (batch 512), weight decay $10^{-4}$, 5\% warm-up + cosine; 970--97,000 steps \\
 & (9,700 unless stated); $\tau_u=0.5$, $\tau_\ell=0.1$; \union draws 25\% of each batch from the labeled images \\
projection head & Linear(512, 512, no bias)--BN--ReLU--Linear(512, 128); ResNet-50: 2,048 hidden units \\
augmentation & crop (scale 0.2--1), flip, colour jitter (0.4, 0.4, 0.4, 0.1; $p=0.8$), grayscale ($p=0.2$) \\
\simclr{}$\to$\supcon & second stage: \supcon on the labeled images, lr 0.02, 485 steps (2,425 and 9,700 also tested) \\
STL-10 & as above at 96 px; the max-pooling layer after the first convolution is kept; $N=105{,}000$ (205 steps per ``epoch'') \\
Tiny-ImageNet & as above at 64 px with the max-pooling layer; $N=100{,}000$ (195 steps per ``epoch''); 10,000 labels \\
BYOL & SGD, lr 0.24, weight decay $5\times10^{-4}$; projector and predictor 4096--256; target momentum 0.99$\to$1 (cosine) \\
SimSiam & SGD, lr 0.06, weight decay $5\times10^{-4}$; projector 2,048 with BN on the output; predictor 512 at constant lr \\
cross-entropy & SGD with Nesterov momentum, lr 0.2 (batch 512), weight decay $5\times10^{-4}$; crop with padding (size/8) + flip \\
FixMatch & 64 labeled + 448 unlabeled images per step, threshold 0.95, RandAugment + Cutout, SGD-Nesterov lr 0.03, \\
 & weight decay $5\times10^{-4}$, schedule $\cos(7\pi k/16K)$, EMA 0.999; $K=10{,}347$ steps at the reference budget \\
MixMatch & 64 labeled + 64 unlabeled images, two augmentations, $T=0.5$, $\alpha=0.75$, $\lambda_u=75$ (linear ramp), \\
 & Adam lr 0.002, weight decay $0.02\cdot$lr, EMA 0.999; 51,733 steps at the reference budget \\
linear probes & standardised features; multinomial logistic regression, L-BFGS; $\lambda\in\{10^{-6},\dots,10^{-1}\}$ selected \\
 & on 20\% of the labels available to the probe, then refitted on all of them \\
$k$NN & cosine similarity, $k=200$, weights $\exp(s/0.1)$ \\
\ftfive (default) & linear-probe initialisation, SGD-Nesterov lr 0.01, weight decay $5\times10^{-4}$, batch 128, \\
 & 60 passes over the labeled images, crop with padding + flip \\
\ftfive (tuned) & head: LP or random; lr: 0.003, 0.01, 0.03, 0.1; passes: 30, 60, 120; regulariser: none, label smoothing 0.1, \\
 & mixup ($\alpha=0.2$) (72 recipes); widened grid: lr also 0.3 and passes also 240 (120 recipes), used for \supcon \\
 & and \union and for the held-out arm; at 9,700 steps the oracle search of \supcon and \union uses heads LP and random, \\
 & lr 0.03, 0.1, 0.3, passes 60, 120, 240 and label smoothing or mixup (36 recipes); selected on 1,000 of the \\
 & 5,000 labels (or chosen on the test set of seed 0 for the oracle recipe), retrained on all 5,000 \\
self-distillation & teacher: \ftfive (default); student: same architecture, initialised from the pretrained encoder; \\
 & 9,700 steps on all images, batch 512, lr 0.1, weight decay $5\times10^{-4}$, temperature 1 \\
\bottomrule
\end{tabular}
\end{adjustbox}
\end{table}

\section{Choosing the \union variant on a validation split}
\label{app:pilot}
The two design choices of \union (temperatures and the share of labeled images per batch) were made on a
validation split: 5,000 training images were held out, the methods were trained for 4,850 steps on the
remaining 45,000, and the probes were evaluated on the held-out images (\cref{tab:pilot}, one seed). The
test set was not used. Uniform batches, in which about 10\% of the images are labeled, were the second-best
variant, and we kept it as an ablation. On the test set, with three seeds, it reaches \R{unionunif_f0.1_e100:lp50:pm}\%
and \R{unionunif_f0.1_e200:lp50:pm}\% \lpall at 9,700 and 19,400 steps (the chosen variant:
\R{union_f0.1_e100:lp50:pm}\% and \R{union_f0.1_e200:lp50:pm}\%) and \R{unionunif_f0.1_e100:ft:pm}\% and
\R{unionunif_f0.1_e200:ft:pm}\% \ftfive (\R{union_f0.1_e100:ft:pm}\% and \R{union_f0.1_e200:ft:pm}\%). Under the
oracle-label probe the two variants are close (\R{cmp:unionunif-union:lp50:e100:sdiff}\pt at 9,700 steps,
\R{cmp:unionunif-union:lp50:e100:ci}, equivalent within one point; \R{cmp:unionunif-union:lp50:e200:sdiff}\pt at
19,400), but after fine-tuning the uniform variant is better at both budgets
(\R{cmp:unionunif-union:ft:e100:sdiff}\pt, \R{cmp:unionunif-union:ft:e100:ci}, and
\R{cmp:unionunif-union:ft:e200:sdiff}\pt, \R{cmp:unionunif-union:ft:e200:ci}). The validation comparison scored the
variants with linear probes, which hardly separate them, so it could not anticipate this. A likely reason for the
difference is that uniform batches give about 90\% instead of 75\% of every batch to instance discrimination over
all images, the part of the loss that keeps improving the label-free arm. No conclusion changes: after fine-tuning
the uniform variant stays \R{d:simclr-unionunif:ft:e100}\pt (9,700 steps) and \R{d:simclr-unionunif:ft:e200}\pt
(19,400 steps) below \simclr, while under the oracle-label probe at 9,700 steps it is
\R{d:simclr-unionunif:lp50:e100}\pt above it, so with this variant the oracle-label probe would overstate
\union even more.

The same validation split serves to check the length and learning rate of the \simclr{}$\to$\supcon stage of
\cref{sec:where}, which were chosen on test accuracies. A \simclr encoder was
trained on the 45,000 remaining images for 19,400 steps, continued with \supcon on their 10\% labeled subset (4,500
images) for 485, 2,425 or 9,700 steps with three learning rates, and probed on the 5,000 held-out images
(\cref{tab:valstage}). The best setting is 485 steps with learning rate \R{valstage:best:lr}.

\begin{table}[h]
\centering
\small
\caption{Choosing the \simclr{}$\to$\supcon stage on the validation split (45,000 training images, 4,500 labels, one
seed; accuracy in \% on 5,000 held-out training images; the test set is not used). ``LP (4,500 labels)'' is the
counterpart of \lpfive, ``LP (45,000 labels)'' of \lpall.}
\label{tab:valstage}
% generated by scripts/make_assets.py -- do not edit
\begin{tabular}{rccc}
\toprule
second-stage steps & learning rate & LP (4,500 labels) & LP (45,000 labels) \\
\midrule
no second stage & -- & 85.7 & 87.6 \\
485 & 0.005 & 87.2 & 88.7 \\
485 & 0.02 & 87.0 & 88.4 \\
485 & 0.25 & 84.1 & 86.3 \\
2,425 & 0.005 & 86.4 & 87.9 \\
2,425 & 0.02 & 85.5 & 87.5 \\
2,425 & 0.25 & 84.4 & 86.7 \\
9,700 & 0.005 & 84.9 & 86.6 \\
9,700 & 0.02 & 84.7 & 86.4 \\
9,700 & 0.25 & 82.9 & 85.2 \\
\bottomrule
\end{tabular}

\end{table}

\begin{table}[h]
\centering
\small
\caption{Validation-split comparison used to fix the temperatures and the batch composition of \union.}
\label{tab:pilot}
\begin{adjustbox}{max width=\linewidth}
% generated by scripts/make_assets.py -- do not edit
\begin{tabular}{lccc}
\toprule
configuration (validation split, 4,850 steps) & LP (all labels) & $k$NN & LP (10\% labels) \\
\midrule
\simclr, $\tau=0.5$ & 79.2 & 72.8 & -- \\
\simclr, $\tau=0.2$ & 79.8 & 73.8 & -- \\
\simclr, $\tau=0.1$ & 78.9 & 72.4 & -- \\
\supcon, $\tau=0.1$ & 83.0 & 80.9 & 79.7 \\
\supcon, $\tau=0.5$ & 82.4 & 79.6 & 79.3 \\
\union, $\tau=0.1$ for all anchors & 80.9 & 74.4 & 77.4 \\
\union, $\tau=0.5$ for all anchors & 82.5 & 77.7 & 80.2 \\
\union, $\tau_\ell=0.1$, $\tau_u=0.5$, uniform batches & 83.6 & 80.8 & 81.4 \\
\union, $\tau_\ell=0.1$, $\tau_u=0.5$, 50\% labeled per batch & 84.2 & 82.7 & 80.8 \\
\union, $\tau_\ell=0.1$, $\tau_u=0.5$, 25\% labeled per batch (chosen) & 84.5 & 82.8 & 81.6 \\
\bottomrule
\end{tabular}

\end{adjustbox}
\end{table}

\section{Two pitfalls when training \supcon}
\label{app:pitfalls}
Two failure modes cost us several runs and may save time for others.

\paragraph{Half precision overflows at $\tau=0.1$}
With $\tau=0.1$ each exponentiated logit can reach $e^{1/\tau}\approx 22{,}000$, and the denominator of
\cref{eq:loss} sums about a thousand of them. This is far above the largest float16 number (65,504), so a
loss computed under automatic mixed precision returns infinities. We run the forward pass in mixed
precision and compute the loss in float32.

\paragraph{Batch normalisation on the output of the projection head collapses \supcon}
With a batch-normalisation layer after the last linear layer of the head, followed by $\ell_2$
normalisation, all embeddings converged to one point and the loss stayed at $\log(2B-1)$. The official
\supcon head has batch normalisation in the hidden layer only, and we use it for all contrastive methods.

\section{Pool statistics and constructed pools}
\label{app:pools}

\paragraph{Definitions}
Let $\phi(x)$ be the $\ell_2$-normalised feature of image $x$ in the reference space, $S_c$ the pool images
of class $c$ and $T_c$ the test images of class $c$. Then
\begin{align}
\text{coverage radius}(S) &= \frac{1}{C}\sum_{c=1}^{C}\frac{1}{|T_c|}\sum_{x\in T_c}\;\min_{s\in S_c}\bigl(1-\phi(x)^\top\phi(s)\bigr), \\
\text{spread}(S) &= \frac{1}{C}\sum_{c=1}^{C}\frac{1}{|S_c|(|S_c|-1)}\sum_{s\neq s'\in S_c}\bigl(1-\phi(s)^\top\phi(s')\bigr).
\end{align}
The class means of the encoder features in \cref{sec:geometry} are also used for two neural-collapse
statistics that we report in the repository: the deviation from a simplex equiangular tight frame,
$\frac{1}{C(C-1)}\sum_{c\neq c'}\bigl|\cos(\tilde\mu_c,\tilde\mu_{c'})+\tfrac{1}{C-1}\bigr|$ with
$\tilde\mu_c=\mu_c-\frac1C\sum_{k}\mu_k$, and the coefficient of variation of the norms
$\lVert\tilde\mu_c\rVert$ (both are zero for a simplex ETF \citep{papyan2020neural}).

\paragraph{Reference spaces}
The \simclr reference encoder was trained on all images for 9,700 steps (seed 0) and never saw a label. The
supervised reference encoder was trained with cross-entropy on all 50,000 labels.

\paragraph{Constructed pools}
All pools contain exactly 500 images per class.
\emph{Redundant} ($k=10$ or $50$): the images of a class are visited in the random order that defines the
random subset; each visited image that is not yet taken becomes a seed and is added together with its
$k-1$ most similar untaken images of the same class, until 500 images are taken.
\emph{$k$-center}: starting from a random image, greedily add the image whose cosine distance to its nearest
chosen image is largest \citep{sener2018coreset}.
\emph{Mixtures}: a share of 25, 50 or 75\% of each class is chosen by one of the two constructions and the
rest is drawn at random from the remaining images of the class.
Pools are built in the \simclr reference space. \Cref{tab:coverage} lists their statistics and the
accuracies obtained on them.

\begin{table}[h]
\centering
\small
\caption{Pools of 5,000 images with controlled coverage: statistics of the pool and accuracy of \supcon and
\simclr trained on it for 9,700 steps (CIFAR-10; mean $\pm$ standard deviation over seeds).}
\label{tab:coverage}
\begin{adjustbox}{max width=\linewidth}
% generated by scripts/make_assets.py -- do not edit
\begin{tabular}{lccccccc}
\toprule
 & & \multicolumn{2}{c}{coverage radius $\downarrow$} & spread & \multicolumn{2}{c}{\supcon on the pool} & \simclr on the pool \\
\cmidrule(lr){3-4}\cmidrule(lr){6-7}
pool (5,000 images) & seeds & \simclr space & supervised space & & LP-50k & LP-5k & LP-50k \\
\midrule
random & 5 & 0.082 & 0.023 & 0.352 & 84.1{\scriptsize$\,\pm\,$0.3} & 81.5{\scriptsize$\,\pm\,$0.5} & 77.1{\scriptsize$\,\pm\,$0.4} \\
25% k-center & 3 & 0.083 & 0.022 & 0.389 & 83.7{\scriptsize$\,\pm\,$0.2} & 81.1{\scriptsize$\,\pm\,$0.4} & 75.1{\scriptsize$\,\pm\,$0.5} \\
50% k-center & 3 & 0.085 & 0.022 & 0.411 & 83.7{\scriptsize$\,\pm\,$0.4} & 80.1{\scriptsize$\,\pm\,$0.3} & 72.2{\scriptsize$\,\pm\,$0.3} \\
75% k-center & 3 & 0.090 & 0.021 & 0.422 & 83.3{\scriptsize$\,\pm\,$0.4} & 79.2{\scriptsize$\,\pm\,$0.3} & 70.0{\scriptsize$\,\pm\,$0.5} \\
k-center & 2 & 0.105 & 0.021 & 0.426 & 83.2{\scriptsize$\,\pm\,$0.2} & 76.5{\scriptsize$\,\pm\,$0.7} & 67.9{\scriptsize$\,\pm\,$0.4} \\
25% redundant & 3 & 0.083 & 0.023 & 0.340 & 82.8{\scriptsize$\,\pm\,$0.2} & 80.1{\scriptsize$\,\pm\,$0.5} & 76.9{\scriptsize$\,\pm\,$0.5} \\
50% redundant & 2 & 0.085 & 0.024 & 0.332 & 82.6{\scriptsize$\,\pm\,$0.0} & 78.8{\scriptsize$\,\pm\,$0.1} & 76.8{\scriptsize$\,\pm\,$0.5} \\
75% redundant & 2 & 0.089 & 0.024 & 0.324 & 81.8{\scriptsize$\,\pm\,$0.2} & 77.1{\scriptsize$\,\pm\,$0.4} & 76.8{\scriptsize$\,\pm\,$0.1} \\
redundant (k=10) & 1 & 0.087 & 0.025 & 0.335 & 82.1 & 78.3 & 76.9 \\
redundant (k=50) & 2 & 0.099 & 0.024 & 0.317 & 80.8{\scriptsize$\,\pm\,$0.3} & 73.5{\scriptsize$\,\pm\,$0.9} & 77.2{\scriptsize$\,\pm\,$0.1} \\
\bottomrule
\end{tabular}

\end{adjustbox}
\end{table}

\paragraph{Pool size}
For random pools the spread is flat in the pool size while the coverage radius decreases monotonically from
1\% of the data to all of it, in both reference spaces and in pixel space (\cref{tab:poolsize}). Spread therefore
cannot explain why a larger pool is better; the coverage radius can.

\begin{table}[h]
\centering
\small
\caption{Coverage radius and spread of random class-balanced pools of increasing size (CIFAR-10, mean over
seeds; cosine distances in the two reference spaces and in pixel space).}
\label{tab:poolsize}
\begin{adjustbox}{max width=\linewidth}
% generated by scripts/make_assets.py -- do not edit
\begin{tabular}{rrccccccc}
\toprule
 & & & \multicolumn{2}{c}{\simclr space} & \multicolumn{2}{c}{supervised space} & \multicolumn{2}{c}{pixel space} \\
\cmidrule(lr){4-5}\cmidrule(lr){6-7}\cmidrule(lr){8-9}
random pool & images & seeds & radius & spread & radius & spread & radius & spread \\
\midrule
1\% & 500 & 5 & 0.125 & 0.352 & 0.033 & 0.025 & 0.443 & 0.919 \\
5\% & 2,500 & 5 & 0.092 & 0.352 & 0.025 & 0.026 & 0.378 & 0.921 \\
10\% & 5,000 & 5 & 0.082 & 0.352 & 0.023 & 0.025 & 0.356 & 0.920 \\
20\% & 10,000 & 5 & 0.073 & 0.352 & 0.021 & 0.025 & 0.337 & 0.920 \\
50\% & 25,000 & 5 & 0.063 & 0.352 & 0.018 & 0.025 & 0.313 & 0.921 \\
100\% & 50,000 & 1 & 0.057 & 0.351 & 0.016 & 0.025 & 0.296 & 0.920 \\
\bottomrule
\end{tabular}

\end{adjustbox}
\end{table}

\section{Additional results}
\label{app:extra}

\begin{table}[h]
\centering
\small
\caption{All budgets of \cref{fig:budget} (CIFAR-10, 10\% labels, ResNet-18; accuracy in \%, mean $\pm$
standard deviation over seeds).}
\label{tab:budgets}
\begin{adjustbox}{max width=\linewidth}
% generated by scripts/make_assets.py -- do not edit
\begin{tabular}{rcccccccc}
\toprule
 & \multicolumn{2}{c}{\simclr} & \multicolumn{2}{c}{\simclrsub} & \multicolumn{2}{c}{\supcon} & \multicolumn{2}{c}{\union} \\
\cmidrule(lr){2-3}\cmidrule(lr){4-5}\cmidrule(lr){6-7}\cmidrule(lr){8-9}
steps & LP-50k & LP-5k & LP-50k & LP-5k & LP-50k & LP-5k & LP-50k & LP-5k \\
\midrule
970 & 67.1{\scriptsize$\,\pm\,$0.3} & 62.8{\scriptsize$\,\pm\,$0.4} & -- & -- & 80.5{\scriptsize$\,\pm\,$0.4} & 77.8{\scriptsize$\,\pm\,$0.4} & 77.0{\scriptsize$\,\pm\,$0.5} & 73.7{\scriptsize$\,\pm\,$0.4} \\
2,425 & 75.0{\scriptsize$\,\pm\,$0.1} & 71.4{\scriptsize$\,\pm\,$0.2} & -- & -- & 83.3{\scriptsize$\,\pm\,$0.4} & 81.0{\scriptsize$\,\pm\,$0.3} & 83.5{\scriptsize$\,\pm\,$0.4} & 80.7{\scriptsize$\,\pm\,$0.3} \\
4,850 & 79.6{\scriptsize$\,\pm\,$0.5} & 76.4{\scriptsize$\,\pm\,$0.1} & -- & -- & 83.8{\scriptsize$\,\pm\,$0.6} & 81.2{\scriptsize$\,\pm\,$0.4} & 85.0{\scriptsize$\,\pm\,$0.5} & 82.0{\scriptsize$\,\pm\,$0.2} \\
9,700 & 83.7{\scriptsize$\,\pm\,$0.2} & 81.3{\scriptsize$\,\pm\,$0.2} & 77.1{\scriptsize$\,\pm\,$0.4} & 72.2{\scriptsize$\,\pm\,$0.5} & 84.1{\scriptsize$\,\pm\,$0.3} & 81.5{\scriptsize$\,\pm\,$0.5} & 85.6{\scriptsize$\,\pm\,$0.1} & 80.9{\scriptsize$\,\pm\,$0.5} \\
19,400 & 87.8{\scriptsize$\,\pm\,$0.2} & 85.8{\scriptsize$\,\pm\,$0.2} & 73.7{\scriptsize$\,\pm\,$0.2} & 69.2{\scriptsize$\,\pm\,$0.9} & 83.9{\scriptsize$\,\pm\,$0.3} & 81.5{\scriptsize$\,\pm\,$0.3} & 86.2{\scriptsize$\,\pm\,$0.4} & 80.7{\scriptsize$\,\pm\,$0.6} \\
38,800 & 89.9{\scriptsize$\,\pm\,$0.3} & 88.4{\scriptsize$\,\pm\,$0.1} & 72.5{\scriptsize$\,\pm\,$0.1} & 65.8{\scriptsize$\,\pm\,$0.2} & 84.0{\scriptsize$\,\pm\,$0.2} & 81.5{\scriptsize$\,\pm\,$0.4} & 87.0{\scriptsize$\,\pm\,$0.2} & 81.0{\scriptsize$\,\pm\,$0.4} \\
97,000 & 90.5 & 89.6 & -- & -- & 84.0 & 81.2 & -- & -- \\
\bottomrule
\end{tabular}

\end{adjustbox}
\end{table}

{\longtabfont
\setlength{\tabcolsep}{3pt}\setlength{\LTcapwidth}{\linewidth}
\begin{longtable}{@{}lccccc@{}}
\caption{Planned comparisons (CIFAR-10 with 10\% labels unless the group names STL-10; ResNet-18). $\Delta$ is the
difference in mean test accuracy with its hierarchical-bootstrap 95\% interval (100,000 replicates over seeds and
test images); $p$-values are Holm-adjusted over the whole table; the last column is the result of the equivalence
test with a margin of $\pm1$ point. $\Delta$ is computed from unrounded means, so it can differ by 0.1 from the
difference of the rounded means in other tables.}
\label{tab:stats}\\
\toprule
comparison & seeds & $\Delta$ (points) & 95\% CI & $p$ (Holm) & equiv.? \\
\midrule
\endfirsthead
\multicolumn{6}{@{}l}{\emph{\tablename~\thetable{} (continued)}} \\
\toprule
comparison & seeds & $\Delta$ (points) & 95\% CI & $p$ (Holm) & equiv.? \\
\midrule
\endhead
\midrule
\multicolumn{6}{r@{}}{\emph{continued on the next page}} \\
\endfoot
\bottomrule
\endlastfoot
\rawinput generated/tab_stats_rows.tex
\end{longtable}
}

\begin{table}[h]
\centering
\small
\caption{Numbers behind \cref{fig:wins}: accuracy (\%) by label fraction $f$ and budget. \supcon, \union
and fine-tuning use one seed except at $f=10\%$ (three seeds at 2,425 and 38,800 steps, five at 9,700; one for
fine-tuning at 2,425 steps); the \simclr probes use three to five seeds. FT-$f$ was not run with all
labels.}
\label{tab:fraction}
\begin{adjustbox}{max width=\linewidth}
% generated by scripts/make_assets.py -- do not edit
\begin{tabular}{rrccccccccc}
\toprule
 & & \multicolumn{3}{c}{\simclr} & \multicolumn{3}{c}{\supcon} & \multicolumn{3}{c}{\union} \\
\cmidrule(lr){3-5}\cmidrule(lr){6-8}\cmidrule(lr){9-11}
steps & $f$ & LP-all & LP-$f$ & FT-$f$ & LP-all & LP-$f$ & FT-$f$ & LP-all & LP-$f$ & FT-$f$ \\
\midrule
2,425 & 1\% & 75.0 & 64.6 & 64.8 & 69.5 & 55.9 & 55.8 & 76.3 & 62.3 & 59.3 \\
2,425 & 5\% & 75.0 & 68.9 & 74.8 & 80.0 & 75.3 & 76.2 & 80.8 & 76.7 & 77.2 \\
2,425 & 10\% & 75.0 & 71.4 & 79.6 & 83.3 & 81.0 & 82.8 & 83.5 & 80.7 & 81.9 \\
2,425 & 20\% & 75.0 & 72.5 & 84.0 & 86.5 & 85.8 & 86.8 & 84.8 & 83.6 & 85.2 \\
2,425 & 50\% & 75.0 & 74.3 & 89.2 & 89.4 & 88.9 & 90.9 & 84.8 & 84.1 & 89.5 \\
2,425 & 100\% & 75.0 & 75.0 & -- & 89.9 & 89.9 & -- & -- & -- & -- \\
\addlinespace
9,700 & 1\% & 83.7 & 76.2 & 76.6 & 69.3 & 56.2 & 54.5 & 81.7 & 57.5 & 65.4 \\
9,700 & 5\% & 83.7 & 79.9 & 82.4 & 80.0 & 75.7 & 76.7 & 83.8 & 76.5 & 76.4 \\
9,700 & 10\% & 83.7 & 81.3 & 85.7 & 84.1 & 81.5 & 82.4 & 85.6 & 80.9 & 81.9 \\
9,700 & 20\% & 83.7 & 82.1 & 88.7 & 87.6 & 86.5 & 87.9 & 87.9 & 86.2 & 87.2 \\
9,700 & 50\% & 83.7 & 83.2 & 91.9 & 91.4 & 91.2 & 92.7 & 90.9 & 90.4 & 92.1 \\
9,700 & 100\% & 83.7 & 83.7 & -- & 93.5 & 93.5 & -- & -- & -- & -- \\
\addlinespace
38,800 & 1\% & 89.9 & 84.6 & 75.6 & 70.2 & 56.4 & 54.2 & 86.3 & 48.1 & 59.7 \\
38,800 & 5\% & 89.9 & 87.3 & 88.5 & 80.1 & 75.2 & 75.5 & 86.0 & 75.9 & 76.9 \\
38,800 & 10\% & 89.9 & 88.4 & 90.3 & 84.0 & 81.5 & 82.4 & 87.0 & 81.0 & 81.5 \\
38,800 & 20\% & 89.9 & 89.1 & 91.8 & 88.0 & 86.8 & 87.9 & 88.7 & 85.9 & 86.5 \\
38,800 & 50\% & 89.9 & 89.7 & 93.8 & 92.2 & 91.9 & 93.1 & 91.2 & 90.2 & 92.2 \\
38,800 & 100\% & 89.9 & 89.9 & -- & 94.5 & 94.5 & -- & -- & -- & -- \\
\bottomrule
\end{tabular}

\end{adjustbox}
\end{table}

\begin{table}[h]
\centering
\small
\caption{Choosing the fine-tuning recipe (CIFAR-10, 5,000 labels, encoders at 19,400 steps; grid of 72
recipes, widened to 120 recipes, with learning rates up to 0.3 and up to 240 passes, for the oracle choice of
\supcon and \union). ``Label split'': recipe chosen on 1,000 of the 5,000 labels; the selection score is the validation
accuracy. ``Test set (oracle)'': recipe chosen on the test set of seed 0, an upper bound; the selection score
is the test accuracy of seed 0, and the last column reports the other two seeds. ``Recipe of \simclr'': the
label-split choice for \simclr applied to another encoder. Test accuracy is the mean $\pm$ standard deviation
over seeds.}
\label{tab:ftsearch}
\begin{adjustbox}{max width=\linewidth}
% generated by scripts/make_assets.py -- do not edit
\begin{tabular}{lllcc}
\toprule
encoder & recipe chosen on & chosen recipe & range of the selection score over the grid & test accuracy \\
\midrule
\simclr & label split & LP head, lr 0.01, 30 passes, label smoothing & 82.3--89.6 & 89.4{\scriptsize$\,\pm\,$0.1} \\
\simclr & test set (oracle) & LP head, lr 0.01, 120 passes, label smoothing & 85.4--89.8 & -- \\
\addlinespace
\addlinespace
\addlinespace
\supcon & label split & LP head, lr 0.003, 30 passes, none & 93.2--100.0 & 81.3{\scriptsize$\,\pm\,$0.3} \\
\supcon & test set (oracle) & LP head, lr 0.1, 120 passes, mixup & 80.9--87.3 & -- \\
\addlinespace
\union & label split & LP head, lr 0.003, 120 passes, label smoothing & 89.6--100.0 & 82.1{\scriptsize$\,\pm\,$0.1} \\
\addlinespace
\bottomrule
\end{tabular}

\end{adjustbox}
\end{table}

\begin{table}[h]
\centering
\small
\caption{Regularising \supcon at 38,800 steps (CIFAR-10, 10\% labels). The last column is the best $k$NN
accuracy over the checkpoints of the run, an upper bound on early stopping.}
\label{tab:reg}
\begin{adjustbox}{max width=\linewidth}
% generated by scripts/make_assets.py -- do not edit
\begin{tabular}{lccccc}
\toprule
\supcon, 38,800 steps & seeds & LP-50k & LP-5k & $k$NN-50k (final) & $k$NN-50k (best checkpoint) \\
\midrule
default (weight decay $10^{\ensuremath{-}4}$) & 3 & 84.0{\scriptsize$\,\pm\,$0.2} & 81.5{\scriptsize$\,\pm\,$0.4} & 82.4{\scriptsize$\,\pm\,$0.2} & 82.4 \\
weight decay $5\times10^{\ensuremath{-}4}$ & 1 & 84.2 & 81.8 & 82.1 & 82.2 \\
weight decay $10^{\ensuremath{-}3}$ & 1 & 84.4 & 82.2 & 82.8 & 82.9 \\
stronger augmentation & 1 & 83.8 & 81.1 & 82.4 & 82.4 \\
weight decay $5\times10^{\ensuremath{-}4}$ + stronger augmentation & 1 & 83.4 & 80.7 & 81.9 & 81.9 \\
\bottomrule
\end{tabular}

\end{adjustbox}
\end{table}

\begin{table}[h]
\centering
\small
\caption{Steps, architecture and learning rate (CIFAR-10, 10\% labels). 900 steps are ``100 epochs'' over the
5,000 labeled images; 970 steps are ``10 epochs'' over all images.}
\label{tab:r50}
% generated by scripts/make_assets.py -- do not edit
\begin{tabular}{llrcccc}
\toprule
encoder & method & steps & learning rate & seeds & LP-50k & LP-5k \\
\midrule
ResNet-18 & \supcon & 900 & 0.25 & 3 & 79.7{\scriptsize$\,\pm\,$0.8} & 77.2{\scriptsize$\,\pm\,$0.6} \\
ResNet-18 & \supcon & 970 & 0.25 & 3 & 80.5{\scriptsize$\,\pm\,$0.4} & 77.8{\scriptsize$\,\pm\,$0.4} \\
ResNet-18 & \supcon & 9,700 & 0.25 & 5 & 84.1{\scriptsize$\,\pm\,$0.3} & 81.5{\scriptsize$\,\pm\,$0.5} \\
ResNet-18 & \supcon & 19,400 & 0.25 & 5 & 83.9{\scriptsize$\,\pm\,$0.3} & 81.5{\scriptsize$\,\pm\,$0.3} \\
ResNet-18 & \simclr & 970 & 0.25 & 3 & 67.1{\scriptsize$\,\pm\,$0.3} & 62.8{\scriptsize$\,\pm\,$0.4} \\
ResNet-18 & \simclr & 9,700 & 0.25 & 5 & 83.7{\scriptsize$\,\pm\,$0.2} & 81.3{\scriptsize$\,\pm\,$0.2} \\
ResNet-18 & \simclr & 19,400 & 0.25 & 5 & 87.8{\scriptsize$\,\pm\,$0.2} & 85.8{\scriptsize$\,\pm\,$0.2} \\
ResNet-50 & \supcon & 900 & 0.25 & 1 & 73.2 & 68.9 \\
ResNet-50 & \supcon & 900 & 0.05 & 1 & 58.8 & 54.5 \\
ResNet-50 & \supcon & 9,700 & 0.25 & 1 & 86.4 & 82.9 \\
ResNet-50 & \supcon & 19,400 & 0.25 & 2 & 85.7{\scriptsize$\,\pm\,$0.2} & 82.6{\scriptsize$\,\pm\,$0.0} \\
ResNet-50 & \simclr & 970 & 0.25 & 1 & 61.7 & 54.7 \\
ResNet-50 & \simclr & 9,700 & 0.25 & 1 & 85.8 & 82.6 \\
ResNet-50 & \simclr & 19,400 & 0.25 & 2 & 90.0{\scriptsize$\,\pm\,$0.1} & 87.8{\scriptsize$\,\pm\,$0.0} \\
\bottomrule
\end{tabular}

\end{table}

\section{Colour reliance of \simclr}
\label{app:colour}
\Cref{tab:colour} removes the colour augmentations of \simclr one at a time. With crop and flip only, the
encoder reaches \R{simclr_exp2_e100:lp50}\% and loses most of it when the test images are converted to
grayscale (\R{simclr_exp2_e100:linear_acc_gray}\%); its features are closer to the colour histogram of the
image than to the class label (linear CKA \R{simclr_exp2_e100:cka_hist} against
\R{simclr_exp2_e100:cka_label} \citep{kornblith2019cka}). A linear classifier on the $8^3$-bin colour
histogram alone reaches \R{histprobe}\%. Random grayscale removes most of the dependence on colour, and
colour jitter and grayscale together give the default accuracy. These runs use one seed and confirm earlier
reports \citep{chen2020simclr,chen2021intriguing,robinson2021can,geirhos2020shortcut}.

\begin{table}[h]
\centering
\small
\caption{Colour reliance of \simclr under four augmentation pipelines (CIFAR-10, 9,700 steps). ``Gray test''
evaluates the \lpall probe on grayscale test images. CKA is linear centred kernel alignment between the
features and the colour histogram or the one-hot label; histogram $R^2$ is the variance of the histogram
explained by a ridge regression on the features.}
\label{tab:colour}
\begin{adjustbox}{max width=\linewidth}
% generated by scripts/make_assets.py -- do not edit
\begin{tabular}{lcccccc}
\toprule
 & & \multicolumn{2}{c}{LP-50k accuracy} & \multicolumn{3}{c}{colour reliance} \\
\cmidrule(lr){3-4}\cmidrule(lr){5-7}
\simclr augmentation & seeds & colour test & gray test & CKA(histogram) & CKA(label) & histogram $R^2$ \\
\midrule
crop + flip & 1 & 60.9 & 24.5 & 0.57 & 0.16 & 0.64 \\
crop + flip + colour jitter & 1 & 75.9 & 38.2 & 0.47 & 0.28 & 0.48 \\
crop + flip + grayscale & 3 & 77.3{\scriptsize$\,\pm\,$0.4} & 75.8{\scriptsize$\,\pm\,$0.2} & 0.31 & 0.34 & 0.33 \\
crop + flip + colour jitter + grayscale (default) & 5 & 83.7{\scriptsize$\,\pm\,$0.2} & 82.6{\scriptsize$\,\pm\,$0.2} & 0.22 & 0.47 & 0.25 \\
\bottomrule
\end{tabular}

\end{adjustbox}
\end{table}

\section{How published comparisons state their budgets}
\label{app:audit}
To see whether unequal compute is a problem in practice, we read the experimental sections of 26 papers
(2019--2025) in which one arm is trained on a labeled subset, or with a labeled subset, and is compared with
an arm trained on a larger pool. Ten papers turned out not to make such a comparison; among them is
\citet{cole2022when}, whose pools of different size are all used for self-supervised pretraining without
labels. \Cref{tab:audit} classifies the other 16. The reading was done from the text of the papers only,
without their code, by one reader; ``unclear'' means that we could not recover the budget of at least one arm
from the paper, not that the comparison is unfair. The one paper in the category of unmatched comparisons
was read a second time in full (arXiv version 1), which confirmed the verdict: SupCon is trained on the 1\% of
the data that is labeled and SimCLR on 70\%, both for 50 epochs; the paper discusses the difference in the number
of images but not in gradient steps, and it does not state the batch size. The quotations behind every verdict,
with their location, are in the repository.

\begin{table}[h]
\centering
\small
\caption{Budget reporting in 16 published comparisons between a labeled-subset arm and a larger-pool arm.}
\label{tab:audit}
\begin{adjustbox}{max width=\linewidth}
\begin{tabular}{p{0.20\linewidth} p{0.46\linewidth} p{0.30\linewidth}}
\toprule
verdict & papers & typical statement of the budget \\
\midrule
matched or accounted for (7) &
\citet{assran2020suncet}, \citet{assran2021paws}, \citet{zheng2022simmatch}, \citet{fini2023semi},
\citet{newell2020useful}, \citet{liu2022self}, \citet{wang2022usb} &
iterations equalised, cost per iteration equalised, or separate columns for labeled and unlabeled epochs \\
\addlinespace
unequal but disclosed (3) &
\citet{zhai2019s4l}, \citet{Cai2022SemiViT}, \citet{zoph2020rethinking} &
epochs defined on the labeled set, or each arm trained to convergence \\
\addlinespace
equal epochs over pools of different size (1) &
\citet{Benato2023Linking} &
the same number of epochs for every pool, no discussion of steps \\
\addlinespace
unclear (5) &
\citet{chen2020simclrv2}, \citet{henaff2020data}, \citet{li2021comatch}, \citet{berthelot2019mixmatch},
\citet{Zhang2021CoCalibration} &
one ``epochs'' column without a definition, or no budget for the baseline \\
\bottomrule
\end{tabular}
\end{adjustbox}
\end{table}

The trap described in \cref{sec:compute} is therefore not common in this sample: most well-known papers
match iterations or report the two budgets separately. The more frequent issue is that the unit of the
budget is not stated in a way that lets a reader check the match. Our sample is small and leans towards
widely cited self- and semi-supervised papers, so it says little about application papers, where
\citet{Benato2023Linking} is an example of the unmatched comparison.
